\section{DIF-SASRec Pretraining: Iterative Development}
\label{sec:dif_sasrec_training}

This section documents the iterative pretraining process for the DIF-SASRec sequential recommendation model.
Three distinct training runs were conducted, each informed by a root-cause analysis of the previous run.
The progression illustrates the importance of the capacity-to-data ratio in transformer pretraining and motivates the final 100k-user training configuration adopted by the system.

\subsection{Motivation}

The prior recommendation engine (GRU-SeqDQN) was an online reinforcement-learning model that required live user traffic to learn meaningful item associations.
In the cold-start state — before any interactions were recorded — it produced effectively random recommendations, making offline benchmarking impossible and delaying personalisation for every new deployment.

DIF-SASRec replaces this with an offline-pretrained transformer that ships warm from day one.
The model is trained on historical Amazon Reviews 2023 Books interactions using a sampled softmax loss, then further fine-tuned per user at inference time via the AgentPool mechanism (Section~\ref{sec:agent_pool}).

\subsection{Training Run 1: Initial Baseline}
\label{sec:train_run1}

The first run established a working baseline using a lightweight 2-block architecture trained on 20{,}000 users.
Users were selected from the Books 5-core subset of the Amazon Reviews 2023 dataset~\cite{hou2024blair}, filtered to items present in the FAISS index, and split by leave-one-out: each user's last interaction is held out as the test item, the preceding sequence is used for training.

\begin{table}[H]
\centering
\caption{DIF-SASRec Run~1: configuration and result.}
\label{tab:train_run1}
\begin{tabular}{ll}
\hline
\textbf{Parameter}       & \textbf{Value} \\
\hline
Parameters               & 2{,}273{,}075 \\
Hidden dimension         & 256 \\
Transformer blocks / heads & 2 / 4 \\
Training users           & 20{,}000 (avg.\ 40--50 clicks) \\
Batch size               & 512 \\
Negatives per step       & 256 (fixed pool) \\
Optimiser                & Adam, lr = 1e-3 (fixed) \\
Epochs                   & 30 \\
\hline
\textbf{HR@10}           & \textbf{0.2973} \\
\textbf{NDCG@10}         & \textbf{0.1668} \\
\hline
\end{tabular}
\end{table}

At HR@10 = 0.2973, DIF-SASRec was 3$\times$ the random baseline (0.1000), confirming that sequential modelling was functioning.
However, it remained well below the content-similarity baseline (HR@10 = 0.4522), suggesting that the 2-block, 256-dimensional architecture was under-capacity for the task.

\subsection{Training Run 2: Scale-Up Attempt and Regression}
\label{sec:train_run2}

Based on the Run~1 result, the model was scaled up: hidden dimension doubled from 256 to 512, transformer blocks increased from 2 to 4, and the training set expanded to 200{,}000 users.
Additional training improvements were also applied simultaneously: the optimiser was changed from Adam to AdamW with weight decay 0.01, a linear warmup with cosine decay LR schedule was introduced, and Automatic Mixed Precision (AMP) was enabled.
Table~\ref{tab:train_run2} compares the two configurations.

\begin{table}[H]
\centering
\caption{DIF-SASRec Run~1 vs.\ Run~2: configuration comparison.}
\label{tab:train_run2}
\begin{tabular}{lll}
\hline
\textbf{Parameter}   & \textbf{Run~1} & \textbf{Run~2} \\
\hline
Parameters           & 2.3M           & 12.4M \\
Hidden dimension     & 256            & 512 \\
Blocks / heads       & 2 / 4          & 4 / 8 \\
Negatives            & 256            & 512 \\
Optimiser            & Adam           & AdamW (wd=0.01) \\
LR schedule          & Fixed 1e-3     & Cosine (2-epoch warmup) \\
Batch size           & 512            & 1{,}024 \\
Training users       & 20{,}000       & 200{,}000 \\
Avg.\ clicks         & 40--50         & 21 \\
Final training loss  & 3.35           & 3.06 \\
\hline
\textbf{HR@10}       & 0.2973         & \textbf{0.1950} \\
\hline
\end{tabular}
\end{table}

The result was a regression: HR@10 dropped from 0.2973 to 0.1950 despite lower training loss and larger architecture.

\textbf{Root cause analysis.}
Two compounding factors were identified.

\emph{Overfitting due to capacity/data mismatch.}
A misconfiguration in the data setup script left the evaluation user list at 20{,}000 entries — the same training population as Run~1.
A model with 5$\times$ more parameters trained on an identical dataset overfits: training loss improved (3.35 $\to$ 3.06), but held-out generalisation to the last-item test degraded.

\emph{Interaction density below signal floor.}
When the data selection was corrected to truly include 200{,}000 users, average clicks per user dropped from 40--50 to approximately 21.
A 4-block transformer with \texttt{MAX\_SEQ\_LEN} = 50 requires sufficiently long input sequences to propagate gradient signal through all layers.
At an average of 21 clicks, a large fraction of users had sequences shorter than the receptive field of the deeper blocks, effectively starving those blocks of training signal.

\subsection{Training Run 3: Calibrated 100k Configuration}
\label{sec:train_run3}

The key insight from Run~2 was that \emph{interaction density matters more than user count} for this architecture.
The data selection cap was adjusted to 100{,}000 users (the top-100k most active readers, sorted by number of rated items).
At this cap, avg.\ train clicks = 33.5 — comfortably above the signal floor while still providing 5$\times$ more user diversity than Run~1.

Three additional training improvements were applied:

\begin{enumerate}
    \item \textbf{Per-epoch negative resampling.}
    Runs~1 and~2 used a fixed pool of 50{,}000 negative items sampled once before training.
    The model gradually memorised which items were easy negatives, causing the training signal to stagnate.
    In Run~3, a fresh 50{,}000-item pool is drawn from the embedding cache at the start of each epoch, continuously restoring the difficulty of the negative task.

    \item \textbf{Batch size 2{,}048 with AMP.}
    Automatic Mixed Precision (fp16 forward pass, fp32 gradient accumulation) halves activation memory, allowing batch size 2{,}048 — doubling GPU throughput compared to Run~2.
    AMP was present in Run~2 but a scheduler-ordering bug suppressed its benefit.

    \item \textbf{AMP + LR scheduler ordering fix.}
    PyTorch's \texttt{GradScaler.step()} silently skips the optimiser update when inf/nan gradients are detected.
    In Run~2, the LR scheduler was still advancing on skipped steps, decoupling the LR trajectory from the actual gradient update count.
    The fix compares the gradient scale before and after \texttt{scaler.step()}: the scheduler advances only when no skip occurred.
\end{enumerate}

\begin{table}[H]
\centering
\caption{DIF-SASRec Run~3: final configuration and result.}
\label{tab:train_run3}
\begin{tabular}{ll}
\hline
\textbf{Parameter}       & \textbf{Value} \\
\hline
Parameters               & 12{,}420{,}345 \\
Training users           & 100{,}000 (avg.\ 33.5 clicks) \\
Batch size               & 2{,}048 \\
AMP                      & On (fp16 forward, fp32 gradients) \\
Negatives                & 512 (resampled each epoch) \\
Optimiser                & AdamW (lr=1e-3, wd=0.01) \\
LR schedule              & Linear warmup 2 epochs $\to$ cosine decay \\
Epochs                   & 30 \\
\hline
\textbf{HR@10}           & \textbf{0.7745} \\
\textbf{NDCG@10}         & \textbf{0.5029} \\
\hline
\end{tabular}
\end{table}

HR@10 reached 0.7745 --- a 161\% improvement over Run~1 and 78\% above the content-similarity baseline (0.4346).
The three-run progression is summarised in Table~\ref{tab:train_summary}.

\begin{table}[H]
\centering
\caption{DIF-SASRec pretraining trajectory.}
\label{tab:train_summary}
\begin{tabular}{p{2.8cm} p{1.8cm} p{1.6cm} p{2.4cm} p{1.6cm} p{1.8cm}}
\hline
\textbf{Run} & \textbf{Params} & \textbf{Users} & \textbf{Avg.\ clicks} & \textbf{HR@10} & \textbf{NDCG@10} \\
\hline
Run~1 (Apr~11) & 2.3M  & 20k  & 40--50 & 0.2973 & 0.1668 \\
Run~2 (Apr~12) & 12.4M & 200k & 21     & 0.1950 & 0.0972 \\
Run~3 (Apr~12) & 12.4M & 100k & 33.5   & \textbf{0.7745} & \textbf{0.5029} \\
\hline
\end{tabular}
\end{table}

\subsection{Training Loss vs.\ Benchmark Quality}

A superficial reading of the final training losses — 3.35 (Run~1), 3.06 (Run~2), 5.34 (Run~3) — would suggest that Run~3 was the weakest of the three.
In fact it is the strongest by a wide margin (HR@10 = 0.7745).
Three separate factors explain why the raw loss numbers are misleading.

\textbf{Different random baselines.}
The training loss is sampled softmax over $K+1$ items (1 positive + $K$ negatives).
For a uniformly random model, the expected loss is $\ln(K+1)$.
Runs~1 used $K=256$ negatives (random baseline $= \ln(257) \approx 5.55$), while Runs~2 and~3 used $K=512$ negatives (random baseline $= \ln(513) \approx 6.24$).
Comparing absolute losses across different $K$ values therefore conflates model quality with task difficulty.
Table~\ref{tab:loss_normalised} reports each run's loss \emph{as a fraction of its own random baseline}, a scale-free measure of how much the model improved over chance.

\begin{table}[H]
\centering
\caption{Training loss normalised by the random baseline $\ln(K+1)$ for each run.}
\label{tab:loss_normalised}
\begin{tabular}{lrrrr}
\hline
\textbf{Run} & \textbf{Negatives $K$} & \textbf{Random baseline} & \textbf{Final loss} & \textbf{Loss / baseline} \\
\hline
Run~1 & 256 & $\ln(257) \approx 5.55$ & 3.35 & 60.4\% \\
Run~2 & 512 & $\ln(513) \approx 6.24$ & 3.06 & 49.0\% \\
Run~3 & 512 & $\ln(513) \approx 6.24$ & 5.34 & 85.6\% \\
\hline
\end{tabular}
\end{table}

By this normalised measure, Run~2 appears to have the best discrimination (49\% of random), yet it produced the worst benchmark.

\textbf{Overfitting inflates apparent training quality.}
Run~2 trained a 12.4M-parameter model on only 20{,}000 unique users — the same population as Run~1.
With 5$\times$ more parameters and an identical data distribution, the model was able to partially memorise the training sequences.
A model that has memorised its training data will show a very low training loss while generalising poorly to held-out users, which is precisely what Run~2's HR@10 = 0.1950 reveals.
Run~3's higher normalised loss (85.6\%) is not a sign of weakness — it is the expected signature of a model that is fitting 100{,}000 \emph{diverse} users and cannot take a shortcut by memorisation.

\textbf{Training and evaluation measure different tasks.}
The training loss requires the model to rank the true next item at position~1 among $K+1$ candidates — an extremely hard task.
The evaluation protocol only requires the item to appear in the top~10 among 100 candidates (1 positive + 99 random negatives).
A model that cannot consistently top-rank among 513 items can still achieve excellent HR@10 when evaluated against only 99 random negatives, because the evaluation pool is both smaller and randomly composed rather than adversarially hard.

Taken together, these three factors explain the seemingly paradoxical result: the run with the highest absolute training loss (Run~3, 5.34) produced the best recommendation quality (HR@10 = 0.7745), while the run with the lowest training loss (Run~2, 3.06) produced the worst (HR@10 = 0.1950).
The lesson is that training loss should always be interpreted relative to the random baseline for the specific negative sampling configuration, and that a generously low loss in a low-data regime is more likely to signal overfitting than genuine capability.

\subsection{Final Checkpoint}

The Run~3 checkpoint trained for 30 full epochs (approximately 9.3 hours on a single GPU) is used as the pretrained starting point for all per-user fine-tuning at inference time.
Periodic checkpoints were saved every 5 epochs, providing safe fallback points for early-stopping analysis.
The final checkpoint file (\texttt{data/dif\_sasrec\_pretrained.pt}) contains 12{,}420{,}345 parameters in the 512-dimensional, 4-block DIF-SASRec architecture described in Section~\ref{sec:infra_improvements}.
